%======================================================================
\subsection{BW-SIM-06 phase 2 --- Built, 31 Oct: prop wash with actuator disks}\label{sec:sim06-wash-built}
%======================================================================

Phase 2 of \S\ref{sec:sim06-plan} is in (commit \code{SIM-06: prop wash with actuator disks}): the study type
\emph{Prop wash}, next to External flow in the Airflow group of New Study. The case writer, the thrust from the
balance and the results are in \file{src/Mod/BruceWorx/bruceworx\_cfd\_wash.py} (no FreeCAD needed to test it),
built on the phase-1 writer, runner, log reader, convergence rule and the cache folder of \S\ref{sec:sim06-built}.
The study frame (\file{bruceworx\_sim.py}, \file{bruceworx\_sim\_gui.py}) and the Airflow panel
(\file{bruceworx\_cfd\_gui.py}) gained the type through a few lines each. Your answers O3 and O4 were followed:
thrust from the \eref{ROB-04} table at hover, a typed thrust as the override, swirl off (see \emph{Open}).

\paragraph{What you can do now.}
\begin{itemize}[nosep]
  \item \textbf{New Study $\to$ Airflow $\to$ Prop wash} on an assembly whose motors are marked with a
        \code{ThrustAxis} (the Mass \& Balance marking, \eref{ROB-03}). Settings: the propeller diameter (one
        diameter for all disks), the speed (\emph{0 = hover}, otherwise a free stream along $+X$ with the angle of
        attack), mesh preset, turbulence, air temperature and altitude; and under Run the thrust per motor (blank
        = the hover share), the disk offset (blank = just above the motor part) and the wake-side domain length.
  \item \textbf{Thrust.} Each motor's disk carries its \emph{hover share}: the weight from the mass report
        split over the motors by the CG offset, exactly what the Mass \& Balance panel shows. The thrust table is what turns that thrust
        into a throttle and a current; the pre-check warns when a share is above a table's maximum
        (\emph{cannot hover}). \emph{Thrust per motor} replaces all shares with one number.
  \item \textbf{Disks.} A disk is a cylinder of cells of the propeller's diameter, three cells thick, centred on
        the motor's axis a tenth of a diameter above the top of the motor part (or the typed offset from the
        motor centre). Its cells are a \code{cellZone} made by \code{topoSet}; the zone carries a constant
        momentum source $-T\hat{\mathbf n}/\rho$ spread over its volume, so the total force is exactly $T$
        whatever the mesh. The mesh is refined to $D/10$, $D/16$ or $D/24$ across each disk (coarse, medium,
        fine) in a cylinder from half a diameter above to 1.5\,D below, and to one level coarser in a wake
        cylinder of radius $D$ down to $6D$.
  \item \textbf{Domain.} A box with open faces in hover (\code{totalPressure} with
        \code{pressureInletOutletVelocity} on all six faces, so air is drawn in and leaves wherever it wants);
        2 characteristic lengths $L$ (the largest size of the model) round it and 4\,$L$ on the wake side.
        In forward flight the phase-1 free-stream box (inlet, outlet, slip sides), $3L$ upstream and $6L$
        downstream, with the disks turned with the model by the angle of attack.
  \item \textbf{Pre-check.} Phase 1's checks, plus a line per disk (thrust, induced velocity $v_i$ and the far
        wake $2v_i$ from momentum theory), the cell count with the share in the disks, a warning if a disk's
        centre is inside a part, and an error when the diameter is blank or no motor is marked.
  \item \textbf{Results.} A force table per part (pressure plus viscous: $F_x$, $F_y$, $F_z$), the
        \emph{download} on each part as a share of the total thrust (the force along the wake, against the thrust
        axis; positive means thrust lost), and the mean velocity along the wake in every disk next to the value
        momentum theory gives. The live plot shows the force along the thrust axis per part. The convergence
        verdict uses the plan's three tests; in hover the third is the net flow through all six faces against
        the flow the disks move. \emph{Report} adds a Prop wash section (setup, convergence, propellers, forces).
\end{itemize}

\paragraph{Benchmark: one disk in free air} (hover, $T=1$\,N, $D=76$\,mm; test
\code{TestBruceWorxCfdWash.TestDiskBenchmark}, registered as validation case \code{B-CF4}). An axisymmetric
$5^\circ$ wedge, 12\,500 cells, 1500 iterations, serial: 40\,s. The reference is momentum theory,
$v_i=\sqrt{T/2\rho A}=9.567$\,m/s at the disk and $2v_i=19.13$\,m/s in the fully contracted wake.
\begin{center}\small
\begin{tabular}{@{}lrrrrr@{}}
\toprule Mesh (cells across $D$) & Cells & Disk $v$ (m/s) & Error & Wake core $v$ (m/s) & Error\\ \midrule
24 & 5\,000 & 10.81 & $+13\,\%$ & 16.46 & $-14\,\%$\\
48 (the test) & 12\,500 & 10.35 & $+8.2\,\%$ & 17.26 & $-9.6\,\%$\\
72 & 21\,600 & 10.22 & $+6.8\,\%$ & 17.32 & $-9.3\,\%$\\
\bottomrule
\end{tabular}
\end{center}
The disk value is the volume-weighted mean axial velocity in the disk's cells; the wake value is the mean over
$r<0.6R$ between 0.9 and 1.1\,D downstream (the axis peak is 18.4\,m/s at $x/D=1.1$, $-3.9\,\%$). The test accepts
10\,\% at the disk and 15\,\% in the wake (your brief); the plan's 5\,\% and 10\,\% are not met on any of the three
meshes. The disk value falls with the mesh and seems to level off near $+6\,\%$, and the wake core sits about
$9\,\%$ low. I read that as turbulent mixing (the $k$--$\omega$ SST shear layer erodes the core before the wake has
contracted fully, and the disk passes more air than the ideal uniform wake would), but did not separate it: a steady
laminar run of the same case does not converge. A thrust of $T/4$ gives half the speed (test
\code{test\_thrust\_scales\_the\_velocity\_as\_the\_square\_root}), which checks the source's normalisation.

\paragraph{Two frames, four disks each} (coarse preset, one core, hover; run through the study as built).
\begin{center}\small
\begin{tabular}{@{}p{0.27\linewidth}rrrrl@{}}
\toprule Frame & Cells & Time & Thrust & Download & Disks $v$ (momentum theory)\\ \midrule
Small X-frame (plate, four 65\,mm arms; test) & 92\,k & 68\,s & 4$\times$0.5\,N & 0.076\,N, 3.8\,\% & 7.37\,m/s (6.77)\\
ROB-03 frame (80\,mm plate, 110\,mm arms, 28\,mm motors; 388\,g) & 242\,k & 572\,s & 4$\times$0.952\,N & 0.798\,N, 21.0\,\% & 9.46\,m/s (9.33)\\
\bottomrule
\end{tabular}
\end{center}
The small frame converged in 456 iterations, its four arms load alike to 0.1\,\% (symmetric frame), and each arm
takes 2.3\,\% of the thrust. The ROB-03 frame did \emph{not} meet the plan's tests in 1500 iterations (residuals
$2\times10^{-3}$ on $U$, forces moving by more than 0.5\,\% on some parts), although the total download had been
steady to 0.2\,\% since iteration 200 and the net flow through the box was 0.002\,\% of the disks' flow. Of the 21\,\%,
each motor takes 3.3\,\% (0.124\,N) and each arm 2.0\,\% (0.075\,N), the plate and battery nothing. The motors are
directly under the disks: the pressure just below a disk is about $T/A-\tfrac12\rho v_i^2\approx157$\,Pa, and a
28\,mm cylinder 7.6\,mm below it feels about 0.1\,N of that. A real propeller has a hub over the motor and moves no air
there, so this uniform-disk number is an upper bound for the motors (see \emph{Open}).

\paragraph{Decided while building.}
\begin{itemize}[nosep]
  \item \textbf{Not \code{actuationDiskSource}.} The plan named that fvOption. In v2606 it is a \emph{wind turbine}
        model: it takes the upstream speed $U$ and the curves $C_p$, $C_t$ and computes $T=2\rho A U^2a(1-a)$ as a
        \emph{sink} of momentum. In hover $U=0$, so it gives no thrust at all, and a propeller has a given thrust,
        not a curve. Each disk is therefore a \code{vectorSemiImplicitSource} on a cell zone, with the thrust as a
        constant force on the air ($-T$ along the motor's \code{ThrustAxis}). \code{radialActuationDiskSource} has the
        same limit; \code{rotorDiskSource} is blade-element theory and needs the blade data.
  \item \textbf{Wedge normalisation.} In a wedge the source acts on $\alpha/360^\circ$ of the disk, so the benchmark
        passes the thrust times $5/360$; the first run, without it, gave 87\,m/s ($\sqrt{72}$ too much).
  \item \textbf{Open faces in hover.} A free stream does not exist, so the box is open on all faces. The mass-flow
        test of phase 1 (inlet against outlet) cannot hold there; the net flow through all six faces is compared
        with the flow the disks move ($v_iA$ summed), and the net flow was 0.00003\,\% (small frame) and 0.002\,\% (ROB-03) of it.
  \item \textbf{Where the disk sits.} The balance knows the motor's centre, not its propeller. The disk is placed at
        the top of the smallest part that contains the motor's centre plus $0.1D$ (7.6\,mm for 76\,mm), unless
        \emph{Disk offset} is typed; a propeller already modelled as a part makes the pre-check warn when the disk
        centre is inside a part.
  \item \textbf{Thrust is the hover share.} The thrust table cannot give the thrust by itself (it maps throttle to
        thrust and needs a thrust to start from). The share of the weight is the thrust the motor must give to hold the
        drone; the table says whether it can. Climbing or manoeuvring thrust is a typed thrust.
  \item \textbf{Code placement.} All new logic is in \file{bruceworx\_cfd\_wash.py}; \file{bruceworx\_cfd.py} only
        dispatches (\code{is\_wash}) so a parallel change to the runner does not collide with it. The validation
        hook registers the \code{cfd} runner on first use (\code{bruceworx\_validation.runner\_for}).
\end{itemize}

\paragraph{Tests} (\code{TestBruceWorxCfdWash}, 30 with one GUI test, 29 headless; the validation hook is commit \code{SIM-06: register the actuator-disk benchmark as validation case B-CF4}): momentum theory (the 9.57 and 19.1\,m/s of chapter 67, forward
flight and climb solved from $T=2\rho Av_i|\mathbf V+v_i\hat{\mathbf w}|$), the disk geometry and its rotation with the
angle of attack, the case writer (dictionaries parse and \code{foamDictionary} reads each one, per-disk force, zone and
probe, snappy cylinders, open faces against free stream, the wake-side margin, presets, pre-check messages), the
benchmark and its $\sqrt T$ scaling, and on the FreeCAD side the study on the ROB-03 X-frame (type registered, hover
default, thrust equals the balance's shares and sums to $mg$, typed overrides, errors for no diameter or no motor,
case written, report section). The GUI test checks the New Study fields (diameter, thrust, offset appear; angle and area hide; speed 0 and back to 10) and the wash table. The four-disk run (68\,s) is skipped unless \code{BW\_CFD\_SLOW=1}. Also run: \code{TestBruceWorxSimCore} and \code{TestBruceWorxSimulation} (their study-type lists now include Prop wash), \code{TestBruceWorxValidation.TestCases} and \code{B-CF4} through \code{run\_set} (PASS, 35\,s).

\paragraph{Open.}
\begin{itemize}[nosep]
  \item \textbf{Swirl} (the tangential force from the motor torque, off by default in your answer) is not built: the
        thrust table has no torque column, and a swirl force needs the cell's position about the axis (a coded
        source). Give me a torque source (the table, or $Q=P/\omega$ from rpm) and it is one more source term.
  \item \textbf{Propeller diameter} is one number in the study. A per-motor diameter would be a second property on
        the marked motor part; say if the SP-01 ever mixes propeller sizes.
  \item \textbf{Run time.} The ROB-03 frame is 242\,k cells at the coarse preset and 572\,s on one core for 1500 iterations; the
        parallel run is another agent's work, and the disks add nothing to it.
  \item \textbf{Convergence of the jet.} A steady RANS wake keeps a little unsteadiness: the ROB-03 run stalls at
        residuals of $2\times10^{-3}$. Options if you want the verdict green: more iterations, a second-order
        scheme only after the first 500, or a time-averaged (unsteady) run. I would first look at whether the
        overlapping arm and motor solids of that test model, which are separate parts, are the cause.
  \item \textbf{Hub.} Add a \emph{hub diameter} so the disk is an annulus (\code{cylinderToCell} has \code{innerRadius}):
        the thrust then acts on $\pi(D^2-d^2)/4$ and the motor under the hub is spared the pressure under the disk. It
        changes the download on every real drone; I did not add it without asking.
  \item \textbf{Results in 3D} (cut plane, streamlines) belong to the results-display work of phase 1.
\end{itemize}

%----------------------------------------------------------------------
\paragraph{Follow-up, 31 Oct evening: hub, per-motor propeller, why the ROB-03 run does not converge, pressure scale.}
%----------------------------------------------------------------------
Your answers to the \emph{Open} list above, in the order of that list (swirl stays off, streamlines are ParaView's job; neither was touched).
\begin{itemize}[nosep]
  \item \textbf{Hub (commit \code{SIM-06: propeller hub}).} A disk is now a ring: \code{cylinderToCell} gets \code{innerRadius}, the
        thrust is spread over the ring's cells, so it acts on $A=\pi(D^2-d^2)/4$. The hub diameter $d$ is the motor part's
        \code{HubDiameter} (mm), else the study's \emph{Hub diameter}, else $0.2\,D$. I chose $0.2\,D$ because the hub of a
        3~inch quad propeller is 12 to 16\,mm on 76\,mm (16 to 21\,\%); a hub of 0.2\,D takes 4\,\% of the area, so $v_i$ rises 2\,\%.
  \item \textbf{Benchmark B-CF4 with the hub.} The reference changes only through the area: $v_i=\sqrt{T/2\rho A_{ring}}$
        $=9.765$\,m/s ($9.567$ without hub) at the ring and $2v_i=19.53$\,m/s in the far wake. The run (same 12\,500-cell wedge,
        $T=1$\,N, $d=15.2$\,mm) gives $10.61$\,m/s at the ring ($+8.7\,\%$) and $17.00$\,m/s in the wake core ($-13\,\%$; the hole
        fills in about one diameter downstream, the on-axis speed is $16.4$ against $17.3$ at $r=0.3R$). The hubless case is
        unchanged ($10.35$, $+8.2\,\%$). New test \code{test\_disk\_with\_a\_hub\_against\_the\_ring\_area\_theory} accepts 10\,\% at
        the ring and 15\,\% in the wake; the wake margin is thin (13\,\%).
  \item \textbf{Propeller diameter per motor.} Nothing existed (only the study's \code{PropDiameter}). \code{mark\_motor} now adds
        \code{PropDiameter} and \code{HubDiameter} (mm, 0 = take the study's value) to the motor part next to \code{ThrustAxis};
        the disk, its offset above the motor and its hub follow each motor's own values, so mixed propeller sizes work. The
        study's diameter may stay blank when every motor part carries its own.
  \item \textbf{ROB-03 frame with the hub} (coarse preset, 3 cores, 1500 iterations, hover, $4\times0.952$\,N): total download
        $0.755$\,N $=19.8\,\%$ of the thrust (before: $0.798$\,N, $21.0\,\%$). Each motor takes $0.111$\,N, $2.9\,\%$ (before
        $0.124$\,N, $3.3\,\%$), each arm $0.076$\,N, $2.0\,\%$ (before $2.0\,\%$), plate and battery $0.1\,\%$. The hub is smaller than
        the 28\,mm motor, so most of the pressure under the disk still reaches the motor; the saving is 12\,\% of the motor load.
  \item \textbf{Convergence of that run: the geometry is not the cause.} The test model's solids do overlap (arms in the
        plate and in each other at the centre; the motors sit on the arm's top face). I ran the same case on (a) the ten separate
        parts and (b) one fused solid (\code{multiFuse}, one watertight body, 241\,628 against 241\,606 cells). checkMesh is
        the same (non-orthogonality max $64.65^\circ$, skewness $3.2$, Mesh OK), the download is the same ($19.87$ against
        $19.84\,\%$) and neither converges; $U_x$ residuals are $2\times10^{-3}$ (separate) and $9\times10^{-3}$ (fused).
        What fails is a swing of the forces of 0.5 to 3.6\,\% of a part's own load over the last 200 iterations (arm and motor 3:
        3.6\,\%; the tiny plate and battery forces, $0.003$\,N, swing by 10 to 30\,\%, which alone fails the plan's rule) while
        the total is steady to 0.2\,\%. This is the unsteadiness of four impinging jets in steady RANS.
  \item \textbf{Solver-side tries} (same case, 3 cores, 4 minutes each): SIMPLE without \code{consistent} and relaxation
        $p\,0.3$, $U\,0.7$: no change ($U_x$ $1.8\times10^{-3}$, swings as before). \code{div(phi,U)} \code{upwind}: residuals drop to
        $4\times10^{-4}$, arm swings to $0.02\,\%$ and motor swings to $0.6\,\%$, but the disks run at $10.5$ against $9.53$\,m/s
        (+10\,\%; linearUpwind gives +5\,\%) and the download is $20.6\,\%$. \code{limitedLinearV 1}: swings as with
        linearUpwind, download $19.4\,\%$. Upwind with 3000 iterations: residuals stay at $4.5\times10^{-4}$, motor swings $0.5$ to
        $0.66\,\%$. None meets the plan's $10^{-4}$ and $0.5\,\%$, so \textbf{I changed no default}: first-order upwind
        trades accuracy for calm, and the verdict would still read ``not converged''.
  \item \textbf{Pressure colour scale (commit \code{SIM-06: pressure colour scale}).} All parts now share one range, the run's own
        minimum and maximum (before, each part fitted itself). The results panel has \emph{Fix pressure scale}, a min and a max in
        Pa and a \emph{This run} button that copies the range just shown; the three values are saved on the study, so
        external flow and prop wash runs are compared with the same colours. It applies at the next \emph{Show in 3D}. FEM's colour
        bar has no range setting reachable from Python, so each saved mm copy of a patch has $p$ clamped to the range and two
        points that no cell uses, carrying the end values; the surface is unchanged and the values in the case are untouched.
\end{itemize}

\paragraph{Tests} added: ring area and induced velocity, hub through \code{Disk} dict and rotation, \code{innerRadius} in
\code{topoSetDict}, the ring benchmark, motor-part diameter and hub (study, part, default), the New Study hub field,
pressure range over patches, pinned copy (range exact, one polygon, two anchor points), \code{pressure\_scale}, the
fixed range in the 3D view, and the panel fields remembered per study.

\paragraph{Open.}
\begin{itemize}[nosep]
  \item The convergence verdict of a prop-wash study reads ``not converged'' on every real frame. Options: judge the forces on the
        total and on parts above, say, 2\,\% of the thrust (the battery's 0.003\,N is noise); and for the residuals, accept
        $10^{-3}$ for jets, or average the last 500 iterations. Say which; I did not weaken the plan's test myself.
  \item A transient (PIMPLE) run of the frame would show whether the wake is truly unsteady; it is a larger job.
  \item The wake-core check of the ring benchmark is 13\,\% low: the plan's 10\,\% is not met there either.
\end{itemize}

\paragraph{Your answers (1 Nov).} Convergence of a prop-wash study is judged on the total and on the parts that carry more
than 2\,\% of the thrust, each steady to 0.5\,\% over the last iterations, with residuals below $10^{-3}$; the verdict reads
``steady average of an unsteady flow'' when the forces pass and the flow still swings. The ring benchmark's wake check is
15\,\%. The hub default stays $0.2\,D$.


\paragraph{Follow-up, 1 Nov.} Items 1 and 2 of your answers are built. The verdict of a prop-wash study
(\code{wash\_convergence} in \file{bruceworx\_cfd\_wash.py}; external flow keeps its own rule and limit $10^{-4}$) is now judged
like this. The forces are taken along the thrust axis, on the total and on every part that carries more than 2\,\% of the total
thrust (its mean over the window, as a share of the thrust). The window is the last 200 iterations (half the run when it is
shorter, never under 50). A force is \emph{steady} when the mean of the window differs from the mean of the 200 iterations before
it by no more than 0.5\,\%; its \emph{swing} is the largest minus the smallest value in the window, as a share of the mean.
The residuals must be below $10^{-3}$, and the net mass flow through the box below 0.5\,\% of the disks' flow. Three outcomes:
\emph{converged} (all of that, and every swing under 0.5\,\%); \emph{steady average of an unsteady flow} (the means are steady
and the flow balances, but a residual stays above $10^{-3}$ or a force swings; the mean and the swing of each judged force are
shown); \emph{not converged} (a mean still drifts, or the flow does not balance). Parts under 2\,\% are listed with their numbers
and do not decide. The results panel shows the verdict word, the mean (and swing) of each judged force and the list of
parts left out; the report has the same in a table (mean, share of the thrust, swing, drift, judged yes or no) under the
verdict, and the setup tables now say $10^{-3}$.
\begin{itemize}[nosep]
  \item \textbf{ROB-03 frame, re-checked on its stored results} (no new run was needed: the comparison case of 31 Oct holds the
        force history; coarse preset, 3 cores, 217\,139 cells, hub $0.2\,D$, $4\times0.952$\,N, 1500 iterations). Verdict:
        \textbf{steady average of an unsteady flow}. Window: iterations 1300 to 1500. $U_x$ residual $2.1\times10^{-3}$ (over the
        limit; $p$ $8\times10^{-5}$). Total download mean $0.7552$\,N, $19.84\,\%$ of the thrust, swing $0.75\,\%$, means
        steady to $0.000\,\%$; mass flow $0.001\,\%$. Judged: total and six parts (motors 0 to 3, arms 0 and 2); the worst
        mean change $0.44\,\%$ (motor 2), the worst swings arm 2 $3.7\,\%$ and motor 2 $2.7\,\%$, the other judged parts
        $0.15$ to $0.4\,\%$. Listed, not judged: arms 1 and 3 ($1.98$ and $1.97\,\%$ of the thrust, just under the line; swings $1.2$
        and $0.2\,\%$), plate ($0.08\,\%$, swing $11\,\%$), battery ($0.03\,\%$, swing $23\,\%$). The same check on the five solver
        variants of that session (consistent off, upwind, linear limited, fused solid, 3000 iterations) reads ``unsteady'' for
        every one (worst mean change $0.22\,\%$). A quick no-hub X frame (95\,515 cells, 4 cores, 33\,s, stopped by its residuals
        at 429 iterations) reads \emph{converged}: total $0.145$\,N ($3.8\,\%$), swing $0.003\,\%$, residuals under $10^{-5}$.
  \item \textbf{Why the mean is compared with the window before, not with its own first half.} The first try (second half of the
        window against the first) made arm 2 fail at $0.57\,\%$ although its 50-iteration block means stay within
        $0.5\,\%$ of each other and its mean moves by $0.05\,\%$ against the 200 before: a swinging signal read in half windows
        reports the phase of the swing, not a drift. Runs shorter than 400 iterations fall back to the halves.
  \item \textbf{Ring benchmark.} The wake-core check was already 15\,\% in the test
        (\code{test\_disk\_with\_a\_hub\_against\_the\_ring\_area\_theory}) and in the \code{v\_wake} reference of B-CF4
        (\file{Resources/Validation/cases.json}); the 10\,\% check at the ring stays. I corrected the basis text of B-CF4 (the core
        sits 10 to 13\,\% low at 1\,$D$, the check allows 15\,\%).
\end{itemize}
Tests added (\code{TestWashVerdict}, plain Python on synthetic force and log files): converged; a small part that swings does not
fail the verdict; residuals above the limit give the unsteady verdict with the right mean; a swinging judged part with steady mean
gives it too, with a 6\,\% swing measured; a drifting mean is not converged; the report table lists the part under 2\,\%; the
panel shows the word, the means and the parts left out.

\paragraph{Your answers (1 Nov, second round).} The line for judged parts drops to 1.5\,\% of the thrust, so mirror-image parts
(arms at 1.97 to 2.0\,\%) are judged together; the plate and the battery stay out. Built: one constant (\code{SIGNIFICANT\_SHARE} $=0.015$) in \file{bruceworx\_cfd\_wash.py}. On the ROB-03 frame nine forces are judged now (total and eight parts, arms 1 and 3 included); the verdict stays \emph{steady average of an unsteady flow}, the worst swing is arm 2 at $3.7\,\%$ (motor 2 $2.7\,\%$, arm 1 $1.2\,\%$, arm 3 $0.2\,\%$) and the worst mean change is motor 2 at $0.44\,\%$; only the plate and the battery are listed.
